\chapter*{Bibliografía}
\addcontentsline{toc}{chapter}{Bibliografía}
\markboth{BIBLIOGRAFÍA}{}

\bibliographystylelib{es-alpha}
\bibliographystylelink{es-alpha}

% \bibliographylib/\bibliographylink (vía multibib) reutilizan la
% "thebibliography" de la clase book, que por defecto abre un
% \chapter*{...} nuevo por cada una. Como aquí queremos que
% "Libros y artículos" y "Enlaces de Internet" sean secciones dentro de
% este único capítulo "Bibliografía", degradamos \chapter a \section (y el
% nivel de \addcontentsline que generan internamente) solo en este ámbito.
\makeatletter
\begingroup
  \let\chapter\section
  \def\@bib@chapterlvl{chapter}
  \let\@bib@stdaddcontentsline\addcontentsline
  \renewcommand{\addcontentsline}[3]{%
    \def\@bib@lvl{#2}%
    \ifx\@bib@lvl\@bib@chapterlvl
      \@bib@stdaddcontentsline{#1}{section}{#3}%
    \else
      \@bib@stdaddcontentsline{#1}{#2}{#3}%
    \fi}
  \bibliographylib{referencias}
  \bibliographylink{referencias}
\endgroup
\makeatother
